\documentclass[10pt,a4paper]{amsart}
\usepackage[margin=19mm]{geometry}
\usepackage{amsmath,amssymb,mathtools}
\usepackage[scr=rsfs,cal=euler]{mathalfa}
\usepackage{xfrac}
\usepackage[unicode,hidelinks,bookmarks=false]{hyperref}
\newcommand{\ie}{i.e.\ }
\newcommand{\cf}{cf.\ }
\newcommand{\resp}{respectively\ }
\newcommand{\Subsection}{\subsection}
\providecommand{\nobreakdash}{\nobreak\hbox{-}}
\setlength{\emergencystretch}{2em}
\multlinegap0pt
\usepackage{bm}
\usepackage{bbm}
\usepackage{dsfont}
\usepackage{xspace}
\usepackage{cancel}
\usepackage{mathtools}
\usepackage{amsthm}
\makeatletter
\@ifundefined{esempio}{\newtheorem{esempio}{Esempio}}{}
\@ifundefined{teorema}{\newtheorem{teorema}{Teorema}}{}
\@ifundefined{lemma}{\newtheorem{lemma}[teorema]{Lemma}}{}
\@ifundefined{oss}{\newtheorem{oss}[esempio]{Remark}}{}
\makeatother
\DeclareMathOperator{\R}{\mathbb{R}}
\newcommand{\abs}[1]{{\left|#1\right|}}
\newcommand{\dH}{d\mathcal{H}^{n-1}}
\title[Talenti comparison results for solutions to $p$-Laplace equa]{Talenti comparison results for solutions to $p$-Laplace equation on multiply connected domains}
\author{Luca Barbato, Francesco Salerno}
\date{Source: arXiv:2504.06103v2 (CC BY 4.0)}
\begin{document}
\maketitle

\noindent\emph{Source and licence.} Talenti comparison results for solutions to $p$-Laplace equation on multiply connected domains. Version arXiv:2504.06103v2 (CC BY 4.0), distributed under the Creative Commons Attribution 4.0 International licence, as recorded in the arXiv licence metadata for this exact version and shown on the abstract page https://arxiv.org/abs/2504.06103v2. Full rights notice: licenses/arxiv-2504.06103-CC-BY-4.0.txt.\par\smallskip

\noindent\textit{Editorial scope.} Counted unit: the Lemma whose source label is \texttt{compsymm} in the article (source0.tex lines 427--436, source label \texttt{compsymm}), delivered together with the article's own complete original proof of it (source0.tex lines 437--492, one \texttt{proof} environment of 56 lines). No mathematics is written, repaired or re-derived: statement and proof are the article's own text line for line. Required context retained uncounted: prob (lines 153--159); probsymm (lines 161--167); brothers1 (lines 273--275); Broz (lines 277--282); distribuzioemu (lines 295--297); distribuzionephi (lines 299--302); eqdebole (lines 375--377); brotherziemer (lines 423--426). Every definition, display and label of those lines is retained unchanged. Omitted from this package and not redistributed: 1--152, 160--160, 168--272, 276--276, 283--294, 298--298, 303--374, 378--422, 493--1067. Nothing inside a retained excerpt is omitted or altered: every retained body is delivered exactly as the article's lines are written, so no omission receipt of any kind accompanies this package. Citations: the retained excerpts carry no citation command at all, so no bibliography entry of the article is transcribed and no reference is redistributed. Reference adaptation: none was needed and none was made; every reference inside the delivered unit resolves inside it, so no unresolved reference or citation is produced. Printed slips kept literal: no printed word, symbol, hypothesis or number of the counted statement or of its proof was altered. Layout and class: the article's own class is \texttt{article} with packages including graphicx, geometry, fontenc, inputenc, babel, amssymb, lipsum, lmodern, amsmath, amsthm, bm, mathtools, braket, appendix; the wrapper is the lane's pinned \texttt{amsart} editorial wrapper at 10pt on A4, which restates the theorem-like environments the retained text needs and adds the article's own macro definitions that the retained text needs (\texttt{\textbackslash R}, \texttt{\textbackslash abs}, \texttt{\textbackslash dH}), with extra wrapper packages (bm, bbm, dsfont, xspace, cancel, mathtools). The delivered numbering is the unit's own. Assets: no retained span contains a figure environment, a graphics-inclusion command, a PDF-page inclusion, a table or a TikZ picture; no image file is copied into this package, the package declares no asset dependency and no image file is redistributed with it. Licence: version arXiv:2504.06103v2 of this article is distributed under the Creative Commons Attribution 4.0 International licence, as recorded in the arXiv licence metadata for this exact version and shown on the abstract page https://arxiv.org/abs/2504.06103v2. A full rights notice is delivered at licenses/arxiv-2504.06103-CC-BY-4.0.txt. Characters: every retained excerpt is pure ASCII, so the delivered engine, XeLaTeX with its default Latin Modern text font, reports no missing character.
\begin{equation}\label{prob}
    \begin{cases}
        -\Delta_p u=f &\text{ in } \Omega\\
        \abs{\nabla u}^{p-2}\frac{\partial u}{\partial\nu}+\beta \abs{u}^{p-2}u=0 &\text{ on }\partial\Omega_0 \\
        u=c_i&\text{ on }\partial\Omega_i
    \end{cases},
\end{equation}

\begin{equation}\label{probsymm}
    \begin{cases}
        -\Delta_p v=f^\sharp &\text{ in } A_\Omega=\Omega_0^\sharp\setminus S^\sharp\\
        \abs{\nabla v}^{p-2}\frac{\partial v}{\partial\nu}+\beta \abs{v}^{p-2}v=0 &\text{ on }\partial\Omega_0^\sharp \\
        v=\overline c &\text{ on }\partial S^\sharp
    \end{cases}
\end{equation}

    \begin{equation}\label{brothers1}
        \infty>-\mu'(t)\geq\displaystyle \int_{u=t}\dfrac{1}{|\nabla u|}d\mathcal{H}^{n-1}
    \end{equation}

\begin{lemma}\label{Broz}
Let $u\in W^{1,p}(\R^n)$, with $p\in(1,+\infty)$. The distribution function $\mu$ of $u$ is absolutely continuous if and only if 
\begin{equation*}
    \abs{\{0<u<  ||u^\sharp||_\infty\}\cap \{|\nabla u^\sharp| =0\}}=0.
\end{equation*}
\end{lemma}

\begin{equation}\label{distribuzioemu}
    \mu(t)=\abs{U_t},\quad P_u(t)=Per(U_t),
\end{equation}

\begin{equation}\label{distribuzionephi}
    V_t=\{x\in\Omega^\sharp\,:\,|\tilde v(x)|>t\},\quad
    \phi(t)=\abs{V_t},\quad P_v(t)=Per(V_t).
\end{equation}

    \begin{equation}\label{eqdebole}
        \int_{\Omega_0}\abs{\nabla u}^{p-2}\nabla u\cdot\nabla v\,dx+\beta\int_{\partial\Omega_0}\abs{u}^{p-2}uv\,\dH=\int_{\Omega_0} fv\,dx\quad\forall\,v\in\mathcal{X}.
    \end{equation}

\begin{oss}\label{brotherziemer}
    From the symmetry of the problem \eqref{probsymm} follows that $v$ is a symmetric radially decreasing function. In particular, this implies that the maximum is attained on $\partial S^\sharp$. Then the only critical points are those on $\partial S^\sharp$ and,
    by Lemma \ref{Broz}, the distribution function of $v$ is absolutely continuous. 
\end{oss}

\begin{lemma}
    Let $u$ and $v$ be the solutions to \eqref{prob}-\eqref{probsymm}, respectively. Let $\mu$ and $\phi$ be defined as in \eqref{distribuzioemu}-\eqref{distribuzionephi}. Then for a.e. $t>0$ we have
    \begin{equation}\label{compsymm}
        \gamma_n\phi(t)^{\left(1-\frac{1}{n}\right)\frac{p}{p-1}}=\left(\int_0^{\phi(t)}f^\ast(s)\,ds\right)^\frac{1}{p-1}\left(- \phi'(t)+\frac{1}{\beta^\frac{1}{p-1}}\int_{\partial V_t\cap\Omega_0^\sharp}\frac{1}{\tilde v}\,d\mathcal{H}^{n-1}\right),
    \end{equation}
    while for almost all $t>0$ it holds
    \begin{equation}\label{comp}
        \gamma_n\mu(t)^{\left(1-\frac{1}{n}\right)\frac{p}{p-1}}\leq\left(\int_0^{\mu(t)}f^\ast(s)\,ds\right)^\frac{1}{p-1}\left(- \mu'(t)+\frac{1}{\beta^\frac{1}{p-1}}\int_{\partial U^{ext}_t}\frac{1}{\tilde u}\,d\mathcal{H}^{n-1}\right).
    \end{equation}
\end{lemma}

\begin{proof}
    Let $t>0$ and $t\neq c_i$ for all $i=1,\dots, m$ and $h>0$ be sufficiently small such that if $c_i<t<c_{i+1}$ then $c_i<t+h<c_{i+1}$ (for such choice of $t$, the corresponding super-level set does not intersect the holes, where the gradient of the constant extension vanishes). Define
    \begin{equation*}
        \varphi_h(x)=\begin{cases}
                0&\text{ if } 0\leq \tilde u<t \\
                \tilde u-t&\text{ if } t\leq \tilde u<t+h\\
                h &\text{ if }\tilde u\geq t+h
                    \end{cases}.
    \end{equation*}
    Using $\varphi_h$ as a test in the formulation \eqref{eqdebole}
    \begin{equation}\label{weakphi}
        \int_{\Omega_0}\abs{\nabla \tilde u}^{p-2}\nabla \tilde u\cdot\nabla\varphi_h\,dx+\beta\int_{\partial\Omega_0}\tilde u^{p-1}\varphi_h\,\dH=\int_{\Omega_0}f\varphi_h\,dx,
    \end{equation}
    by the definition of $\varphi_h$, we get
    \begin{equation*}
        \begin{split}
            \int_{\Omega_0}\abs{\nabla \tilde u}^{p-2}\nabla \tilde u\cdot\nabla\varphi\,dx=& \int_{U_t\setminus U_{t+h}}\abs{\nabla u}^p\,dx,\\
            \int_{\partial\Omega_0}\tilde u^{p-1}\varphi\,\dH=&h\int_{\partial U_{t+h}^{ext}}\tilde u^{p-1}\,\dH+\int_{\partial U_{t}^{ext}\setminus\partial U_{t+h}^{ext}}\tilde u^{p-1}(\tilde u-t)\,\dH,\\ \int_{\Omega_0}f\varphi\,dx=&h\int_{U_{t+h}}f\,dx+\int_{U_t\setminus U_{t+h}}f(\tilde u-t)\,dx
        \end{split}
    \end{equation*}
    thus dividing \eqref{weakphi} by $h$ becomes
    \begin{equation}\label{INVENTO}
    \begin{split}
        \frac{1}{h}\int_{U_t\setminus U_{t+h}}\abs{\nabla \tilde u}^p\,dx+\frac{\beta}{h}\int_{\partial U_{t}^{ext}\setminus\partial U_{t+h}^{ext}}\tilde u^{p-1}(\tilde u-t)\,\dH&+\beta\int_{\partial U_{t+h}^{ext}}\tilde u^{p-1}\,\dH\\
        &=\int_{U_{t+h}}f\,dx+\frac{1}{h}\int_{U_t\setminus U_{t+h}}f(\tilde u-t)\,dx.
    \end{split}
    \end{equation}
    We remark that
    \begin{equation*}
        \frac{\beta}{h}\int_{\partial U_{t}^{ext}\setminus\partial U_{t+h}^{ext}}\tilde u^{p-1}(\tilde u-t)\,\dH<\beta\int_{\partial U_{t}^{ext}\setminus\partial U_{t+h}^{ext}}\tilde u^{p-1}\,\dH,
    \end{equation*}
    and the last, as $h\rightarrow 0$, it becomes the integral of $\tilde u^{p-1}$ on a set whose measure goes to zero. Then, letting $h\rightarrow0$ in \eqref{INVENTO}, we have
    \begin{equation*}
        \int_{U_t}f\,dx=\int_{\partial U_t}g(x)\,\dH,
    \end{equation*}
    where 
    \begin{equation*}
        g(x)=\begin{cases}
            \abs{\nabla \tilde u}^{p-1}&\text{ on }\partial U_t^{int} \\
            \beta\tilde u^{p-1}&\text{ on }\partial U_t^{ext}
        \end{cases}.
    \end{equation*}
    We remark that, from the assumptions on $t$ and $t+h$, the function $g(x)$ is always non-zero.
    By Cauchy-Schwarz inequality and the previous identity
    \begin{equation*}
        \begin{split}
            Per(U_t)&=\left(\int_{\partial U_t}g(x)^\frac{1}{p}\frac{1}{g(x)^\frac{1}{p}}\,\dH\right)\leq\left(\int_{\partial U_t}g(x)\,\dH\right)^\frac{1}{p}\left(\int_{\partial U_t}\frac{1}{g(x)^\frac{1}{p-1}}\,\dH\right)^{1-\frac{1}{p}}\\ &\leq\left(\int_0^{\mu(t)}f^\ast(s)\,ds\right)^\frac{1}{p}\left(-\mu'(t)+\frac{1}{\beta}\int_{\partial U_t^{ext}}\frac{1}{u}\,\dH\right),
        \end{split}
    \end{equation*}
    on the other hand, by the isoperimetric inequality, we find
    \begin{equation*}
        n\omega_n^\frac{1}{n}\mu(t)^{1-\frac{1}{n}}\leq\left(\int_0^{\mu(t)}f^\ast(s)\,ds\right)^\frac{1}{p}\left(-\mu'(t)+\frac{1}{\beta}\int_{\partial U_t^{ext}}\frac{1}{u}\,\dH\right).
    \end{equation*}
     Lastly, if $v$ solves \eqref{probsymm}, then \eqref{brothers1} holds with equality sign thanks to
    remark \ref{brotherziemer}. In particular, this implies that for $v$ the last inequality holds as an equality.
\end{proof}

\end{document}
